\documentclass[10pt,a4paper]{amsart}
\usepackage[margin=19mm]{geometry}
\usepackage{amsmath,amssymb,mathtools}
\usepackage[scr=rsfs,cal=euler]{mathalfa}
\usepackage{xfrac}
\usepackage[unicode,hidelinks,bookmarks=false]{hyperref}
\newcommand{\ie}{i.e.\ }
\newcommand{\cf}{cf.\ }
\newcommand{\resp}{respectively\ }
\newcommand{\Subsection}{\subsection}
\providecommand{\nobreakdash}{\nobreak\hbox{-}}
\setlength{\emergencystretch}{2em}
\multlinegap0pt
\usepackage{enumerate}
\usepackage{amssymb}
\usepackage{tikz}
\usepackage{mathtools}
\newtheorem{theorem}{Theorem}
\title{This is the title}
\author{K. Mahesh Krishna}
\date{Source: arXiv:2603.01446v2 (CC BY 4.0)}
\begin{document}
\maketitle

\noindent\emph{Source and licence.} This is the title. Version arXiv:2603.01446v2 (CC BY 4.0), distributed by arXiv under the Creative Commons Attribution 4.0 International licence, http://creativecommons.org/licenses/by/4.0/, as recorded in the arXiv licence metadata for this exact version and shown on the abstract page https://arxiv.org/abs/2603.01446v2. Full licence notice: \texttt{licenses/arxiv-2603.01446-CC-BY-4.0.txt}.\par\smallskip

\noindent\textit{Editorial scope.} Counted unit: the article's theorem NTM (source0.tex lines 140--160). The article's complete original proof of it is delivered with it (source0.tex lines 161--210, one \texttt{proof} environment of 50 lines). No mathematics is written, repaired or re-derived: the statement and the proof are the article's own lines, in the article's own notation, and no retained line is substituted or re-flowed.  No further source-local context is needed: every cross-reference of every retained line resolves inside the two delivered spans.  Nothing was omitted from any retained span, so the package carries no omission record.  No retained line contains a citation, so the wrapper carries no bibliography and no bibliography entry.  No retained line contains a figure environment, an image-inclusion command or drawing code, so no image is delivered and the package declares no asset dependency.  Editorial formatting only, in addition: an AMS \texttt{amsart} preamble that restates the article's own declarations and macros used by the retained lines, namely the environments \texttt{theorem} on separate counters, together with the standard packages those lines need.  The retained environments print in this document as Theorem 1 in order of appearance, whereas the article prints the same items with the numbering of its own full document.  The complete native source of the article as distributed in the arXiv e-print for this exact version is bundled unchanged as \texttt{sources/195563/source0.tex}.
\begin{theorem}\label{NTM}
(\textbf{Non-Archimedean Tarski-Maligranda Inequalities})
Let $\mathcal{X}$ be a NALS over a non-Archimedean valued field $\mathbb{K}$. Assume that the characteristic of the field $\mathbb{K}$ is not 2. Then for all $x, y \in \mathcal{X}$, 
\begin{align*}
	\bigg|\|x\|-\|y\|\bigg|&\leq \frac{2}{|2|}\max\{\|x-y\|, \|x+y\|\}-(\|x\|+\|y\|)\\
	&\leq \frac{2}{|2|}\max\{\|x\|, \|y\|\}-(\|x\|+\|y\|)
\end{align*}
and 
\begin{align*}
	\bigg|\|x\|-\|y\|\bigg|\leq \|x\|+\|y\|-\frac{2}{|2|}\bigg|\|x+y\|-\|x-y\|\bigg|.
\end{align*}	
In particular, if $|2|=1$, then 
\begin{align*}
\bigg|\|x\|-\|y\|\bigg|&\leq 2\max\{\|x-y\|, \|x+y\|\}-(\|x\|+\|y\|)\\
&\leq 2\max\{\|x\|, \|y\|\}-(\|x\|+\|y\|)
\end{align*}
and 
\begin{align*}
\bigg|\|x\|-\|y\|\bigg|\leq \|x\|+\|y\|-2\bigg|\|x+y\|-\|x-y\|\bigg|.
\end{align*}	
\end{theorem}

\begin{proof}
	Let $x, y \in \mathcal{X}$. Then 
	\begin{align}\label{F1}
		\|x\|+\|y\|+\bigg|\|x\|-\|y\|\bigg|=2\max\{\|x\|, \|y\|\}.
	\end{align}	
We also have 
\begin{align}\label{A1}
	\max\{\|x-y\|, \|x+y\|\}\geq \|(x+y)+(x-y)\|=\|2x\|=|2|\|x\|
\end{align}
and 
\begin{align}\label{A2}
	\max\{\|x-y\|, \|x+y\|\}\geq \|(x+y)-(x-y)\|=\|2y\|=|2|\|y\|.
\end{align}
Inequalities (\ref{A1}) and (\ref{A2}) give 
\begin{align}\label{F2}
		\max\{\|x-y\|, \|x+y\|\}\geq |2|\max\{\|x\|, \|y\|\}.
\end{align}
Inequalities (\ref{F1}) and (\ref{F2}) give 
\begin{align*}
	\|x\|+\|y\|+\bigg|\|x\|-\|y\|\bigg|=2\max\{\|x\|, \|y\|\}\leq \frac{2}{|2|}\max\{\|x-y\|, \|x+y\|\}.
\end{align*}
Rearranging, 
\begin{align*}
\bigg|\|x\|-\|y\|\bigg|\leq \frac{2}{|2|}\max\{\|x-y\|, \|x+y\|\}-(\|x\|+\|y\|).
\end{align*}
We next see that 
	\begin{align}\label{E1}
	\|x\|+\|y\|-\bigg|\|x\|-\|y\|\bigg|=2\min\{\|x\|, \|y\|\}.
\end{align}	
Further, 
\begin{align}\label{G1}
	\bigg|\|x+y\|-\|x-y\|\bigg|\leq	\|(x+y)+(x-y)\|=\|2x\|=|2|\|x\|
\end{align}
and 
\begin{align}\label{G2}
	\bigg|\|x+y\|-\|x-y\|\bigg|\leq	\|(x+y)-(x-y)\|=\|2y\|=|2|\|y\|.
\end{align}
Inequalities (\ref{G1}) and (\ref{G2}) give 
\begin{align}\label{E2}
	\bigg|\|x+y\|-\|x-y\|\bigg|\leq	|2|\min\{\|x\|, \|y\|\}.
\end{align}
Inequalities (\ref{E1}) and (\ref{E2}) give 
\begin{align*}
		\bigg|\|x+y\|-\|x-y\|\bigg|\leq	|2|\min\{\|x\|, \|y\|\}= \frac{|2|}{2}\left(\|x\|+\|y\|-\bigg|\|x\|-\|y\|\bigg|\right).
\end{align*}
Rearranging, 
\begin{align*}
\bigg|\|x\|-\|y\|\bigg|\leq \|x\|+\|y\|-\frac{2}{|2|}\bigg|\|x+y\|-\|x-y\|\bigg|.
\end{align*}
\end{proof}

\end{document}
